% Chapter 2: Related Work

\chapter{Related Work}\label{ch:related_work}

This chapter reviews the related literature according to the logical layers of this thesis's method, echoing the three classes of challenges raised in \cref{sec:problem_statement}: \cref{sec:ftc} and \cref{sec:pomdp_methods} correspond to ``partial observability,'' introducing the classical frameworks of traditional control theory and partially observable decision-making; \cref{sec:residual_rl}, \cref{sec:robust_rl_dr}, and \cref{sec:il_hil} correspond to ``deployment constraints,'' providing the methodological foundation for the real-robot deployment of this thesis along three threads---the residual structure of robotic RL, robust RL and domain randomization, and imitation learning with human-in-the-loop; \cref{sec:safe_rl} corresponds to ``human--robot collaborative safety,'' providing the background on safe RL and reward hacking, which offers conceptual support for the analysis of the backup-policy reward evolution in \cref{sec:backup_reward}. This chapter focuses on the landscape of the field and the positioning of the method, and does not repeat the specific configuration details of \cref{ch:task_policy}--\cref{ch:real_chapter}.

\section{Fault-Tolerant Robot Control (FDI/FTC)}\label{sec:ftc}

Fault Detection, Isolation, and Reconfiguration (FDI/FTC) is a classical framework for fault-tolerant robot control~\cite{blanke2016diagnosis}. FDI detects faults in a system and locates the fault source through analytical redundancy, state observers, or parameter identification; FTC, once a fault is detected, maintains system stability through control-law reconfiguration. This framework has been applied successfully to problems such as sensor drift and actuator lock-up, but exhibits three key limitations when applied to the ``encoder systematic bias of unknown magnitude'' addressed in this thesis:

\begin{itemize}
  \item \textbf{It requires an explicit fault model}---bias sources addressed in this thesis, such as ``replacement of components without recalibration,'' lack a prior model;
  \item \textbf{It targets ``maintaining stability''}---\emph{classical passive FTC} stops or degrades operation after detection; although active FTC pursues performance retention, it still relies on an explicit fault model, which does not match the ``unknown magnitude with no fault model'' subproblem of this thesis;
  \item \textbf{It has limited ability to detect a latent-variable bias that is constant within an episode}---parameter-identification FDI (such as RLS or UKF) can in principle estimate a constant parameter online, but in the absence of an independent external observation there exists a non-identifiable subspace (consistent with the translational symmetry given in \cref{prop:non_identifiable}).
\end{itemize}

\noindent Haddadin et al.~\cite{haddadin2017robot} survey the problem of robot collision detection and isolation, which likewise focuses on the ``collision event'' rather than the ``residual bias after a collision.'' This thesis adopts a path different from FDI/FTC---it treats the bias as part of the environment distribution and lets the policy adapt online by having seen the bias during training, thereby bypassing the two-stage paradigm of ``detect first, then switch.''

\section{POMDP Solution Methods}\label{sec:pomdp_methods}

Solving partially observable MDPs has long been a core difficulty of reinforcement learning~\cite{sutton2018reinforcement} and sequential decision-making. Kaelbling et al.~\cite{kaelbling1998planning} propose an exact solution method based on belief states, which converts a POMDP into an MDP solved over the belief space; such methods have theoretical guarantees on small-scale discrete problems, and extending them to the high-dimensional continuous state space used in this thesis requires approximation (such as point-based POMDP solvers or deep variational belief models). Hausknecht and Stone~\cite{hausknecht2015deep} propose the Deep Recurrent Q-Network (DRQN), which uses an RNN/LSTM to implicitly memorize the history of observations, achieving good results on partially observable tasks such as Atari. Treating a latent-variable bias that is constant within an episode as a \emph{static hidden parameter} or a \emph{contextual MDP} is another thread for handling this class of problem (it shares with the randomized-bias training of this thesis the basic idea of ``exposing the policy to multiple parameters during training,'' but explicitly modeling the latent variable requires an additional inference module). Weighing engineering trade-offs (implementation complexity, hyperparameter sensitivity, and compatibility with the existing BC/RL training stack), this thesis chooses the path of observation augmentation plus end-to-end RL---introducing an independent external observation so that the POMDP is approximately restored to an MDP in the information-theoretic sense (see \cref{sec:observability} for details); the other paths above are in principle solvable and are mentioned as future work in \cref{ch:conclusion}.

\section{Robotic Reinforcement Learning and Residual Structure}\label{sec:residual_rl}

Residual reinforcement learning (Residual RL) was proposed by Johannink et al.~\cite{johannink2019residual}: a hand-designed base controller first completes most of the action, and RL then learns a residual policy that makes small corrections to the base controller's output. Residual RL can significantly reduce sample complexity and improve training stability on many robotic tasks, and is a common compromise in the sim-to-real field---it leaves the analytically modelable part to the base controller and the hard-to-model disturbances to RL.

The encoder bias of this thesis happens to corrupt the joint readings used by the base controller (the Jacobian and forward kinematics); whether Residual RL remains effective in this scenario requires a separate empirical analysis. This thesis does not conduct a systematic comparison of Residual RL against BC plus HG-DAgger on the real robot, and leaves it as future work (see \cref{ch:conclusion}).

\section{Robust Reinforcement Learning and Domain Randomization}\label{sec:robust_rl_dr}

Domain Randomization (DR) is the mainstream approach for narrowing the sim-to-real gap. Tobin et al.~\cite{tobin2017domain} propose visual DR on an object pose estimation task---by randomizing parameters such as textures, lighting, and camera poses, the policy is made robust to differences in visual appearance; subsequent work extends this to state-space DR (position noise, control delay, dynamics parameters), achieving sim-to-real success in scenarios such as legged robots and drones.

The ``randomized-bias training'' used in this thesis (see \cref{sec:random_bias} for details) is formally cognate with DR---both rely on the basic idea of ``letting the policy see samples of the deployment distribution during training.'' It is necessary, however, to point out an essential information-theoretic difference between the two: in the original setting of Tobin et al. (pose estimation, grasping), the randomized visual parameters are \emph{redundant variables} for the policy (they should be ignored, and what the policy learns should be an invariance); whereas the bias $\vect{b}$ randomized in this thesis is a latent variable on which the policy \emph{must condition}---according to the conditional-compensation requirement of \cref{subsec:identifiability_vs_solvability}, the optimal policy must output different actions for different values of $\vect{b}$. In this setting the two belong to two opposite information-theoretic paradigms, invariance learning and conditional-compensation learning respectively, and are cognate only in the surface idea that ``the training distribution covers the deployment distribution.'' State-space DR is used as a mature approach in the backup-policy training of \cref{sec:backup_training}---low-dimensional state perturbations such as position noise, velocity jitter, and detection delay can be covered by DR alone to bridge the sim-to-real gap, which forms a contrast with the visual path of the task policy that requires BC plus HG-DAgger real-robot imitation.

\section{Imitation Learning and Human-in-the-Loop}\label{sec:il_hil}

Imitation Learning (IL) and the Human-in-the-Loop (HIL) variants of reinforcement learning (RL) together form the methodological foundation of the real-robot deployment of this thesis. This section reviews two parallel branches---BC$\to$DAgger$\to$HG-DAgger in the IL direction and RLPD$\to$HIL-SERL in the RL$+$HIL direction---and finally gives the rationale for the dual-end choice of this thesis.

\paragraph{IL branch: BC, DAgger, HG-DAgger} \textbf{Behavior Cloning} (BC) is the simplest IL method---supervised learning on expert demonstration data that maps states to actions. BC is simple to implement and stable to train, but its behavior is undefined in out-of-distribution (OOD) states (the compounding-error problem), and it usually requires a large number of demonstrations or additional regularization. \textbf{DAgger} (Dataset Aggregation)~\cite{ross2011dagger}, proposed by Ross et al., is a classical improvement---during policy deployment it queries the expert on newly encountered states, merges the expert actions into the dataset, and retrains, iterating until convergence. The key limitation of DAgger is that the query is initiated by the policy (query-by-policy)---frequently asking the expert for actions in unexpected states while the policy is not yet stable places an undue burden on the human expert. \textbf{HG-DAgger} (Human-Gated DAgger)~\cite{kelly2019hgdagger} is a human-initiated variant of DAgger---the intervention trigger is decided by the human (human-initiated): the operator intervenes only on the failure modes of the policy, and the intervention frames are merged into the next retraining round. This thesis adopts HG-DAgger as the final approach for the real-robot stage.

\paragraph{RL+HIL branch: RLPD, HIL-SERL} \textbf{RLPD}~\cite{ball2023rlpd}, building on Soft Actor-Critic~\cite{haarnoja2018sac}, adds a 50/50 mixed sampling of an offline demonstration buffer and an online buffer, providing demo guidance for sparse-reward tasks. \textbf{HIL-SERL}~\cite{luo2025hilserl}, on top of RLPD, makes human-in-the-loop intervention a continuous source of training data throughout the entire online stage, with intervention transitions reinjected into the offline buffer on an equal footing with the original demos. HIL-SERL achieves outstanding results on several real-robot manipulation tasks in the original work of Luo et al. \emph{The interpretation of this thesis is the following}: the SAC online stage of HIL-SERL requires a relatively high real-robot throughput for the critic to calibrate the Q-surface online through continuous online interaction (inferred backward from the real-robot failure mode in \cref{subsec:hilserl_failure}), and this condition is not easily satisfied on a low-throughput real robot.

\paragraph{The dual-end choice of this thesis} The simulation side adopts HIL-SERL (\cref{sec:sim_task_train})---simulation throughput is high and bias injection is controllable, so HIL-SERL fully converges and validates H1--H3; the real-robot side switches to BC plus HG-DAgger (\cref{sec:real_task_deploy})---under the real-robot conditions of this thesis, the actor drifts substantially after being unfrozen in the SAC online stage of HIL-SERL (a detailed post-mortem is given in \cref{subsec:hilserl_failure}), whereas HG-DAgger does not depend on a critic, has a smaller tuning surface, and is markedly more stable in the small-data real-robot scenario. This choice is determined by the three axes of (observation modality $\times$ data scale $\times$ hardware throughput), and is one instance of the decision principle for choosing a paradigm in Contribution 2 of \cref{sec:contributions}.

\section{Safe Reinforcement Learning and the Reward Hacking Background}\label{sec:safe_rl}

Safe reinforcement learning is concerned with guaranteeing that a system does not violate safety constraints during training or deployment. Methods such as the Constrained MDP and Constrained Policy Optimization (CPO) handle constraints explicitly through techniques such as Lagrangian multipliers or rejection sampling. This thesis does not directly adopt the constrained MDP---the safety requirement of this thesis (human--robot collaborative avoidance) is realized through a \emph{hierarchical structure} (task policy $+$ backup policy $+$ a runtime distance-triggered hard switch, see \cref{sec:backup_paradigm} for details) rather than constrained optimization, which better matches the need for engineering interpretability in industrial deployment scenarios.

Kiemel et al.~\cite{kiemel2024safe} propose, in the context of safe RL for industrial robots, the reward-structure principle of \emph{non-negative immediate rewards $+$ a large termination penalty}, used to train backup-as-safety-layer avoidance policies. The V3 backup-policy reward of \cref{sec:backup_reward} directly inherits this principle, and over the three iterations V1$\to$V2$\to$V3 empirically demonstrates that ``if the reward design does not satisfy this principle, the policy will, through reward misspecification, find an optimal solution outside the design intent.''

\textbf{Reward hacking} refers to a policy violating the design intent while maximizing the reward signal---a typical manifestation is that the policy finds a shortcut behavior that ``raises the reward without truly completing the task.'' Reward hacking is a classical problem in reward engineering, and is widely discussed in both large language model alignment and robotic RL. The two failed iterations V1 (a proximity reward) and V2 (an all-negative reward) of \cref{sec:backup_reward} provide two typical manifestations of reward misspecification (V2 additionally exhibits a termination-policy bias); the specific mechanisms and a comparison of training curves are given in \cref{sec:backup_reward}.
